\section{Discussion}
\label{sec:discussion}

\subsection{Why Does Low-Rank Adaptation Underperform on Prithvi?}
Our split-QKV ablation isolates two effects. The \emph{implementation trap} reflects a property of PyTorch's \texttt{nn.MultiheadAttention}, in which query, key, and value weights are packed into a single \texttt{in\_proj\_weight} tensor. A naive application of LoRA to named \texttt{nn.Linear} children therefore touches only \texttt{out\_proj}, leaving the attention projections effectively frozen. This is a silent failure mode: training proceeds, loss decreases, and accuracy settles only slightly above linear probing, which can easily be misread as a genuine negative result.

Once Q, K, and V are explicitly exposed, split-QKV LoRA improves EuroSAT s2\_full from 0.658 to 0.707 but remains below both BitFit (0.702) and Houlsby (0.821). We attribute this residual gap to two factors. First, the masked autoencoding objective~\cite{he2022mae} used by Prithvi may produce features whose task-relevant directions are not well captured by low-rank updates to attention alone. Second, the frozen backbone restricts LoRA to reshaping existing attention patterns, whereas Houlsby adapters introduce non-linear capacity inside each block. This is consistent with the unified view of PEFT proposed in prior work~\cite{he2022unified}, which predicts that methods with both extra capacity and non-linearity can dominate under sufficiently different pre-training distributions. The completed capacity sweep strengthens this interpretation: split-QKV LoRA remains essentially flat at 0.706--0.709 OA from $r=4$ to $r=64$ while its trainable parameter count rises from 302{,}602 to 4{,}726{,}282. Houlsby, in contrast, increases monotonically from 0.864 OA at $d=8$ (164{,}458 parameters) to 0.901 OA at $d=64$ (1{,}197{,}322 parameters). Thus, the largest LoRA model still trails the smallest Houlsby adapter by 15.6 OA points despite using 28.7 times more trainable parameters, supporting an adapter-placement and non-linearity explanation rather than a simple parameter-budget explanation in this tested setting.

\subsection{When Does Input-Stage Adaptation Help?}
GeoAdapter v2 places a small residual module before the patch embedding. Conceptually, this is attractive for modality bridging, but the empirical picture is regime-dependent. On EuroSAT, where the input is 13-band Sentinel-2 close to Prithvi's HLS pretraining modality, GeoAdapter underperforms linear probing on every modality configuration: the input-stage residual has limited new information to add, and the small parameter budget is insufficient to jointly learn modality alignment and useful feature transformation. On the Linhe production corpus (Section~\ref{sec:linhe}), where the input is 3-band RGB padded with zeros into the HLS template --- a substantial modality shift --- GeoAdapter recovers $+0.087$ mIoU over linear probing at only 4{,}756 parameters, second only to Houlsby's bottleneck adapters. We therefore refine the negative result: input-stage adaptation is fragile when the modality shift is small and the patch embedding is already well-tuned, but becomes useful as the input drifts further from the backbone's pretraining distribution. The recommendation is to treat GeoAdapter as a deployment-time decision: profile the modality gap first, then decide. Two underlying mechanisms still bound its potential. First, the module lies on the critical gradient path to a frozen backbone, so any early-stage misadaptation propagates unchallenged through 12 transformer blocks. Second, the adapter has no clear pre-training prior, unlike bottleneck adapters that operate on already well-structured hidden states. Future modality bridging work should combine input-stage and backbone-level adaptation rather than rely on either alone.

\subsection{Architecture-Awareness as a Methodological Principle}
The fused-QKV trap has an uncomfortable implication: published PEFT comparisons on Prithvi-like backbones may silently conflate \emph{method failure} with \emph{implementation failure}. We believe this motivates a broader methodological recommendation. Before reporting negative results for PEFT methods on any new architecture, authors should report (i) exactly which parameter tensors were registered as trainable, (ii) whether the attention projections are fused, and (iii) whether a diagnostic split variant was attempted. Our study adopts this protocol and we encourage its adoption for future GeoFM benchmarks.

\subsection{Practical Guidance and Limitations}
For practitioners adapting Prithvi-100M to heterogeneous inputs, the safest default among the evaluated configurations is Houlsby adapters with bottleneck dimension 64, which reach 94.5\% of the full fine-tuning ceiling on EuroSAT s2\_full while training 1.4\% of parameters in the original benchmark and 0.901 OA in the completed capacity sweep. When spectral richness is available, preserving it yields larger returns than switching between weak PEFT methods: a linear probe on s2\_full outperforms sophisticated adaptation on sar\_only. LoRA should not be used on fused-QKV backbones without explicit Q/K/V exposure and, even then, should be justified empirically against adapter baselines and parameter-capacity controls.

Our study has four limitations. First, we evaluate a single backbone, Prithvi-100M; extension to SatMAE~\cite{cong2022satmae}, Scale-MAE~\cite{reed2023scalemae}, and SpectralGPT~\cite{hong2024spectralgpt} is necessary before the ranking can be generalized to GeoFMs as a class. Second, although we cover classification (single-label on EuroSAT, multi-label on BigEarthNet-S2) and dense prediction (semantic segmentation on LandCover.ai, Linhe, and LoveDA), regression and temporal forecasting remain open. Third, our BigEarthNet-S2 subsample (10K/5K) is chosen for Colab-scale tractability and may under-represent the long tail of rare classes; full-scale replication is an obvious next step. Fourth, the Linhe evaluation uses Esri-derived supervisory labels with a temporal offset relative to the 2025 imagery. This is useful for testing an operational supervision workflow, but it does not replace independent manual validation of land-cover accuracy.
